\scanpage{18}
Dans tout ceci, il reste cependant un
type important de \deleted{\unsure{hom.}} morphismes, entre
surfaces \inserted{top.} \unsure{comp. clos} ou entre schémas en
courbes, qui n'a pas été pris en ligne
de compte --- savoir les morphismes
d'inclusion ouverte
\[
(23)\qquad X=Y\setminus T\subset Y
\]
correspondants à une opération de \uline{bouchage
de trous}. Mais on est confronté, pour élargir la
catégorie en conséquence, par l'impératif de
pouvoir composer les morphismes étales finis,
et d'inclusion ouverte (type ``bouchage de trous'')
à envisager des morphismes $X\longrightarrow Y$ qui
s'insèrent dans un diagramme commutatif
\diagram{\begin{tikzcd}[ampersand replacement=\&]
X\arrow[r,hook]\arrow[d]\&\widehat X\arrow[d]\\
(24)\quad Y\arrow[r,hook]\&\widehat Y
\end{tikzcd}}
avec $\widehat X\longrightarrow\widehat Y$ fini ramifié, $X\mathbin{\unsure{-|\longrightarrow}}Y$ lui-$=$ étale.
Posant $S=\widehat X\setminus X$, $T=\widehat Y-\unsure{T|Y}$, de sorte que
\[
(25)\qquad S\supset\widehat X|T
\]
Les ennuis techniques proviennent du fait que
l'on est obligé d'admettre que $\widehat X/\widehat Y$ puisse avoir des
pts de ramification effective non seulement en
des pts de $\widehat X|T$, mais aussi en des pts de $\unsure{S|C}$ qui
ne sont \uline{pas} dans $\widehat X|T$.
\margin{17}
